% !PRINTED 243
\weberpage{243}
If $x_1$ and $x_2$ are both imaginary, but not conjugate, then the product
\[
(x_1 - x_2)(x_1' - x_2')
\]
and its square are positive.

But if, finally, $x_1$ and $x_2$ are conjugate imaginary, then their difference is \gesperrt{purely imaginary} and its square negative.

Thus in the product by which we have defined the discriminant there occur exactly as many negative factors as there are pairs of conjugate imaginary roots, and from this we conclude the important fundamental theorem:

\gesperrt{If $f(x) = 0$ is a real equation without multiple roots, then the discriminant is positive or negative according as the number of pairs of conjugate imaginary roots is even or odd.}

The case in which only real roots are present belongs to those in which the discriminant is positive.

For the cases of the quadratic and cubic equations, in which only one pair of imaginary roots can occur, this principal theorem already completes the determination among the various cases possible with regard to the reality of the roots.

\begin{center}
\S.\,77.

\gesperrt{Discussion of the quadratic and cubic equation.}
\end{center}

For the quadratic equation
\[
a_0 x^2 + a_1 x + a_2 = 0
\]
we have (\S.\,46)
\[
\begin{array}{l}
D = a_1^2 - 4a_0a_2 > 0 \quad \text{real roots.}\\
\phantom{D = a_1^2 - 4a_0a_2 } < 0 \quad \text{imaginary roots.}
\end{array}
\]
For the cubic equation
\[
a_0 x^3 + a_1 x^2 + a_2 x + a_3 = 0
\]
we have
\[
D = a_1^2a_2^2 + 18a_0a_1a_2a_3 - 4a_0a_2^3 - 4a_1^3a_3 - 27a_0^2a_3^2,
\]
\[
\begin{array}{l}
D > 0 \quad \text{three real roots,}\\
D < 0 \quad \text{one real, two imaginary roots.}
\end{array}
\]
